\documentclass[11pt]{article}

\usepackage[margin=1in]{geometry}
\usepackage{amsmath, amssymb}
\usepackage{booktabs}
\usepackage{hyperref}
\usepackage{array}
\usepackage{mdframed}
\usepackage{enumitem}

\definecolor{notebg}{RGB}{255,248,230}
\definecolor{noteborder}{RGB}{200,150,40}
\definecolor{exbg}{RGB}{235,245,255}
\definecolor{exborder}{RGB}{60,110,180}

\newmdenv[backgroundcolor=notebg, linecolor=noteborder, linewidth=1pt,
          innertopmargin=6pt, innerbottommargin=6pt, innerleftmargin=8pt, innerrightmargin=8pt]{notebox}
\newmdenv[backgroundcolor=exbg, linecolor=exborder, linewidth=1pt,
          innertopmargin=6pt, innerbottommargin=6pt, innerleftmargin=8pt, innerrightmargin=8pt]{casebox}

\newcommand{\eq}[1]{\begin{equation}#1\end{equation}}

\title{Fiscal Policy in General Equilibrium:\\
Model Specification and Worked Cases\\
\large A companion to the Baxter \& King (1993) replication app}
\author{}
\date{}

\begin{document}
\maketitle

\begin{abstract}
\noindent This document is written for undergraduates using the interactive
app at \url{https://baxter-king.streamlit.app/} to explore Baxter \& King
(1993, \textit{American Economic Review}), ``Fiscal Policy in General
Equilibrium.'' Part I lays out the full model the app solves: preferences,
technology (including the Section VI extension for productive public
capital), the government's budget, the household's optimality conditions,
the closed-form steady state, the log-linear transition dynamics, and how
the app's three comparison panels map onto that math. Part II works
through seven cases that are easy to get wrong intuitively --- why a
spending \emph{composition} change is not economically neutral even when
public capital is unproductive; why the \emph{financing rule} moves the
economy even when every dollar of tax revenue is handed straight back; why
several steady-state variables always move by the exact same percentage by
mathematical necessity, not by coincidence; why the app can lose its
saddle path altogether at extreme income-tax rates; why a ``temporary''
spending increase does not bring the economy back to where it started
once the underlying policy regime has changed; why the real interest
rate shapes fiscal multipliers even though it is a completely separate
sidebar control; and why longer-lasting temporary shocks produce
\emph{bigger}, not smaller, effects on impact. Part III derives
every first-order condition and steady-state equation from scratch.
\end{abstract}

\tableofcontents

\part{Model specification}

\section{Notation}

\begin{center}
\begin{tabular}{cl}
\toprule
Symbol & Meaning \\
\midrule
$C_t$ & Consumption \\
$N_t$ & Labor supplied (fraction of the time endowment); $L_t=1-N_t$ is leisure \\
$Y_t$ & Output \\
$K_t$ & Private capital stock \\
$K^G_t$ & Public (government) capital stock \\
$I_t,\ I^G_t$ & Private and public investment \\
$G_{T,t}$ & Lump-sum government transfers \\
$w_t,\ r_t$ & Real wage and real rental rate on capital \\
$\tau_t$ & Tax rate \\
$\theta_N,\ \theta_K=1-\theta_N$ & Labor's and capital's shares of income \\
$\theta_G$ & Productivity of public capital (Section VI extension) \\
$\theta_L$ & Weight on leisure in utility (calibrated, not a free parameter) \\
$\delta$ & Depreciation rate (same for $K$ and $K^G$) \\
$A$ & Total factor productivity, normalized to $1$ \\
\bottomrule
\end{tabular}
\end{center}

A hat over a variable, e.g.\ $\hat y_t$, denotes a percent deviation from
steady state --- the units used in every transition-path chart in the app.

\section{Preferences}

\eq{ u(C_t,N_t) = \ln C_t + \theta_L \ln(1-N_t) }

Households value consumption and leisure, with log utility in each --- the
standard King--Plosser--Rebelo specification that keeps balanced-growth
labor supply well behaved. Notice that government spending does \emph{not}
appear here at all: households get no direct utility from transfers
$G_{T,t}$ or public investment $I^G_t$. This single fact is the seed of
almost everything counterintuitive in Part II --- government spending only
ever affects the household \emph{indirectly}, through the resource
constraint and the budget constraint, never through a warm-glow ``I like
government spending'' term.

\begin{notebox}
\textbf{Why is there no ``basic government purchases'' category?} Many
textbook treatments split $G$ into wasteful ``purchases'' $G_B$, productive
investment $I^G$, and transfers $G_T$. Here $G_B$ is dropped because it
would be economically identical to $I^G$ at $\theta_G=0$: neither enters
utility, and both use up real resources one-for-one. Keeping a separate
$G_B$ category would just be relabeling. Total government spending in this
model is therefore always exactly two pieces: $G = I^G + G_T$.
\end{notebox}

\section{Technology}

\eq{ Y_t = A\,K_t^{\theta_K}\,(K^G_t)^{\theta_G}\,N_t^{\theta_N}, \qquad \theta_K=1-\theta_N }

Output is Cobb--Douglas in private capital and labor, exactly as in the
paper's baseline model. The Section VI extension (and the app's $\theta_G$
slider) adds public capital $K^G_t$ as a third input: it raises the
productivity of every private input, without the government paying for it
again next period --- a non-rival ``public good'' flavor. Setting
$\theta_G=0$ collapses this exactly to the paper's baseline model, where
public investment is a pure resource cost with \emph{no} productivity
effect. Note that $\theta_K+\theta_N=1$ (constant returns to scale in
$(K,N)$ alone) but $\theta_G$ is \emph{extra}: when $\theta_G>0$ the
production function has \emph{increasing} returns to scale in
$(K,K^G,N)$ jointly, which is exactly the source of the model's
``multiplier $>1$'' results in Part II.

\section{Capital accumulation}

\eq{ K_{t+1} = (1-\delta)K_t + I_t \qquad\qquad K^G_{t+1} = (1-\delta)K^G_t + I^G_t }

Both stocks accumulate by the standard perpetual-inventory rule and (for
simplicity) depreciate at the same rate $\delta$.

\section{Resource constraint and the government budget}

\eq{ Y_t = C_t + I_t + I^G_t }

Output is used for private consumption, private investment, and public
investment. \textbf{Transfers do not appear here.} A transfer just moves
resources from the government's books to a household's pocket without using
up any output, so it nets out of the economy-wide resource constraint even
though it is part of the government's budget below.

\eq{ \tau_t Y_t = I^G_t + G_{T,t}, \qquad \tau_t = \frac{I^G_t+G_{T,t}}{Y_t} }

The government's budget balances every period (no debt in this model), and
the tax rate is \emph{defined} as total government spending divided by
output. This holds identically under \emph{both} financing rules below ---
what differs between them is not how much revenue is raised, but whether
raising it distorts the household's decisions.

\eq{ \text{Income tax:}\quad C_t+I_t = (1-\tau_t)\big(w_tN_t+r_tK_t\big) + G_{T,t} }
\eq{ \text{Lump-sum:}\quad C_t+I_t = w_tN_t+r_tK_t - \tau_tY_t + G_{T,t} }

Under \textbf{income tax} financing, $(1-\tau_t)$ is a proportional wedge
that shrinks the \emph{marginal} return to working and saving --- this is
what makes the tax distortionary. Under \textbf{lump-sum} financing, the
same total revenue $\tau_t Y_t$ is instead collected as a poll tax that does
not depend on how much the household works or saves, so the household faces
$w_t$ and $r_t$ directly (no $(1-\tau_t)$ term at all): a pure income
(wealth) effect with no substitution effect at the margin.

\section{Household optimization}

Maximizing lifetime utility subject to the budget constraint gives two
first-order conditions, one per financing rule.

\eq{ \text{Labor-leisure, income tax:}\qquad \theta_L\,\frac{C_t}{1-N_t} = (1-\tau_t)\,w_t }
\eq{ \text{Labor-leisure, lump-sum:}\qquad \theta_L\,\frac{C_t}{1-N_t} = w_t }

The household works until the marginal disutility of an extra hour (left
side) equals its marginal after-tax benefit (right side). Under income tax
financing that benefit is shrunk by $(1-\tau_t)$; under lump-sum financing
there is no such wedge at all.

\eq{ \text{Euler, income tax:}\qquad \frac{1}{C_t} = \beta\,\frac{1+(1-\tau_{t+1})\,\mathrm{MPK}_{t+1}-\delta}{C_{t+1}} }
\eq{ \text{Euler, lump-sum:}\qquad \frac{1}{C_t} = \beta\,\frac{1+\mathrm{MPK}_{t+1}-\delta}{C_{t+1}},\qquad \beta=\frac{1}{1+r} }

There is no expectation operator here: this model has no stochastic
element (no random shocks to technology, policy, or anything else), so
every ``future'' variable is simply its known, perfect-foresight value.
The household is not \emph{guessing} what $C_{t+1}$ or $\mathrm{MPK}_{t+1}$
will be --- under a given policy path it \emph{knows}, and the Euler
equation is an ordinary deterministic difference equation linking today's
consumption to tomorrow's. (This is what makes the transition path in
Section~\ref{sec:steadystate} below solvable by exact saddle-path methods
rather than by simulation.)

Prices are marginal products: $w_t=\theta_N Y_t/N_t$ and
$\mathrm{MPK}_t=\theta_K Y_t/K_t$. In steady state, $(1-\tau)\mathrm{MPK}=r+\delta$
under income tax, and $\mathrm{MPK}=r+\delta$ under lump-sum.

\section{Steady state (closed form)}
\label{sec:steadystate}

Because $\theta_K+\theta_N=1$, dividing the technology equation through by
$N$ leaves the capital/labor ratio $\kappa=K/N$ pinned down by a
capital-Euler condition written in \emph{per-worker} public capital
$K^G/N$ --- but with an extra $N^{\theta_G}$ factor that does \emph{not}
cancel, since $K^G$ is itself funded as a share of \emph{aggregate} output,
not of output per worker:

\eq{ \text{Income tax:}\qquad (1-\tau)\,\theta_K\,A\,(K^G/N)^{\theta_G}\,N^{\theta_G}\,\kappa^{\theta_K-1} = r+\delta }
\eq{ \text{Lump-sum:}\qquad \theta_K\,A\,(K^G/N)^{\theta_G}\,N^{\theta_G}\,\kappa^{\theta_K-1} = r+\delta }

When $\theta_G=0$ the $N^{\theta_G}$ factor is just $1$ and this collapses
to the familiar CRS-in-$(K,N)$ formula, solved for $\kappa$ directly. When
$\theta_G>0$, $K^G=I^G/\delta$ itself depends on output, \emph{and} the
equation above depends on $N$ directly too, so the app solves a small
nested fixed point (first on $\kappa$ given a guess for $N$, then updating
$N$, iterating to convergence).

Everything else then follows in closed form, with
$s_C=1-s_I-s_{I^G}$ the steady-state consumption share (transfers never
reduce $s_C$, since they are not resource-using) and $s_I=\delta\kappa/(Y/N)$
the investment share:

\eq{ \text{Income tax:}\qquad N = \frac{(1-\tau)\,\theta_N}{\theta_L\,s_C + (1-\tau)\,\theta_N} }
\eq{ \text{Lump-sum:}\qquad N = \frac{\theta_N}{\theta_L\,s_C + \theta_N} }
\eq{ Y=\frac{Y}{N}\cdot N,\qquad K=\kappa N,\qquad C=s_C\,Y,\qquad I=s_I\,Y }

$\theta_L$ itself is not a free parameter: it is \emph{calibrated} once, at
a chosen baseline policy, so that steady-state $N$ matches a target
hours-worked share (the paper's Table 1 strategy, and $N_{\text{target}}=20\%$
in the app). Holding that calibrated $\theta_L$ fixed while comparing two
different policies is what lets $N$ move \emph{honestly} in response to a
policy change, rather than being pinned back to the target by construction.

\section{Transition dynamics (log-linearized)}

Around a steady state, every equation above can be approximated to first
order in percent deviations. The production and labor-supply conditions
combine into one static relationship,
\eq{ \hat y_t = y_k\,\hat k_t + y_{kg}\,\hat{k}^G_t + y_c\,\hat c_t + y_g\,\hat g_t, }
capital accumulation becomes a linear law of motion for the
\emph{predetermined} state $\hat k_t$,
\eq{ \hat k_{t+1} = \Phi_{kk}\,\hat k_t + \Phi_{kc}\,\hat c_t + \Phi_{kg}\,\hat g_t + \Phi_{k,kg}\,\hat{k}^G_t, }
and the Euler equation pins down the \emph{forward-looking} (``jump'')
variable $\hat c_t$ in terms of tomorrow's capital and government spending,
\eq{ \hat c_t = \text{coef}_c\,\hat c_{t+1} + \text{coef}_k\,\hat k_{t+1} + \text{coef}_g\,\hat g_{t+1} + \text{coef}_{kg}\,\hat{k}^G_{t+1}. }
Public capital's own law of motion is exogenous, so the system has exactly
one predetermined and one jump variable: it is \textbf{saddle-path stable},
with one stable and one unstable eigenvalue of the reduced-form transition
matrix. The app diagonalizes that matrix and pins the unstable component to
the present value of \emph{future} government-spending shocks (the
King--Plosser--Rebelo / Blanchard--Kahn method) rather than simulating
forward naively, which would blow up numerically.

\section{Reading the app: two comparison panels}
\label{sec:panels}

Every chart in the app is one of the two objects above: the closed-form
steady state of Section~\ref{sec:steadystate}, or the log-linear
transition path of this section. The app presents them as \textbf{three
panels}. The first two differ only in \emph{which two steady states are
being compared} --- the underlying mathematics is identical in both; the
third simply reorganizes Panel 2's own comparison by variable.

\subsection{Panel 1: Cumulative Effects}

Compares the paper's one \emph{fixed} benchmark calibration
($\theta_N=0.58$, $\delta=10\%$, $r=6.5\%$, $G/Y=20\%$, all of it
transfers, lump-sum financing, $\theta_G=0$) against \textbf{whatever the
sidebar is currently set to}. Every control --- financing, composition,
$\theta_G$, $r$, and $G/Y$ --- can differ between the two columns
simultaneously, so this panel reports the \emph{net, cumulative} effect of
everything that has been changed at once. It is the right panel for
questions like Cases A and B below: ``what does switching to income tax
financing do to the economy, holding everything else at the benchmark?''

Its steady-state comparison is the direct application of
Section~\ref{sec:steadystate}: solve for the benchmark's $(\kappa,N,Y,C,\ldots)$
once, solve for the current settings' $(\kappa,N,Y,C,\ldots)$ once
(re-using the benchmark's calibrated $\theta_L$ on both sides, so that
labor supply is free to respond honestly rather than being pinned back to
the same target by construction), and report percent differences.

Its \emph{time-series} chart applies the log-linear transition machinery
of this section in a way that goes beyond a simple $\Delta G/Y$
experiment. Financing, composition, $\theta_G$, and $r$ are always treated
as a \textbf{permanent, immediate switch in regime} --- this model has no
mechanism for a structural parameter to ``revert,'' so there is no
sensible ``temporary $\theta_G$'' or ``temporary financing rule.'' The
economy is imagined to start ($t=0$) with the \emph{benchmark's} capital
stocks $K,K^G$ but the \emph{current settings'} structural parameters and
policy from that point on, with the transition linearized around the
\emph{destination} steady state:
\eq{ \hat k_0 = \ln\!\frac{K_{\text{benchmark}}}{K_{\text{current}}}, \qquad \hat{k}^G_0 = \ln\!\frac{K^G_{\text{benchmark}}}{K^G_{\text{current}}}. }
Only total government spending $G$ respects the sidebar's Duration control
(item 4.2), since a temporary spending change is economically meaningful
even when the surrounding policy regime is not. If Permanent, $\hat g_t=0$
for all $t$ ($G/Y$ jumps to the current share immediately and stays there,
exactly the convention Panel 2 uses for a permanent shock). If Temporary,
$G/Y$ instead reverts to the \emph{benchmark's} $20\%$ share after the
chosen number of years --- see Case D below for why this is more subtle
than it first looks.

\subsection{Panel 2: The Marginal Effects of Government Spendings}

Compares $G/Y=20\%$ against \textbf{the sidebar's current $G/Y$}, holding
financing, composition, $\theta_G$, and $r$ fixed at whatever the sidebar
currently has them at, \emph{on both sides} of the comparison. This
isolates the effect of the $\Delta G/Y$ lever specifically: resetting
$G/Y$ to $20\%$ always brings this panel back to exactly zero, no matter
what the other controls are set to. It is the panel behind the app's
headline multipliers ($\Delta Y/\Delta G$) and the Table-3-style
duration-sensitivity replication.

Its time series is the ordinary $\Delta G/Y$ transition path: $\hat g_t$ is
the (possibly temporary) government-spending shock itself, $\hat k_0=0$
(capital starts exactly at the pre-shock steady state, since only $G/Y$
changed, nothing structural), and the linearization is around whichever
steady state the shock is being compared to (the new steady state, for a
permanent shock; the old one, for a temporary shock).

\subsection{Panel 3: Multiplier Effects}

Takes the \emph{same} pair of steady states as Panel 2 --- $G/Y=20\%$ vs.
the sidebar's current $G/Y$, everything else held fixed --- and reports,
for each of $X\in\{Y,C,I,K,K^G\}$, the multiplier
\eq{ \frac{\Delta X}{\Delta G} = \frac{X_{\text{current}}-X_{20\%}}{G_{\text{current}}-G_{20\%}}, }
both \textbf{long-run} (using the two steady-state levels directly) and
\textbf{on impact} (using each variable's year-0 level along Panel 2's own
transition path in place of $X_{\text{current}}$). No new steady state or
transition path is computed here --- it is purely Panel 2's comparison,
reorganized by variable instead of collapsed into a single headline
number. A multiplier above $+1.00$ for output means a marginal dollar of
spending raises long-run output by \emph{more} than a dollar, the paper's
central result; investment's impact multiplier can be large even when its
long-run multiplier is small, since capital only needs to move once to
reach the new steady state (the ``accelerator boom'').

\begin{notebox}
\textbf{Why not just one panel?} A single ``current settings vs.\ some
baseline'' panel cannot serve every purpose at once. If the baseline
tracks the current settings on everything except $G/Y$ (Panel 2's
convention), then changing financing \emph{alone} --- with $G/Y$ left at
$20\%$ --- shows up as \emph{no change at all}, since both sides of the
comparison would use the same financing rule. If instead the baseline is
the one truly fixed benchmark (Panel 1's convention), resetting $G/Y$ to
$20\%$ will not zero out the panel unless every other control also happens
to match the benchmark. Splitting the two eliminates this ambiguity:
Panel 1 answers ``how far is the economy, in total, from the textbook
case?''; Panel 2 answers ``what, specifically, is this one lever doing, in
aggregate?''; Panel 3 answers the same question as Panel 2, but
``\ldots and how does that break down variable by variable?''
\end{notebox}

\part{Worked cases}

The cases below are the ones students most often get wrong by intuition
alone. Cases A--C follow directly from the closed-form steady state of
Section~\ref{sec:steadystate}; Cases D--E require the transition dynamics
of Section~\ref{sec:panels} as well, and are best explored live in the
app rather than by hand.

\section{Case A: composition is \emph{not} neutral once $I^G>0$, even at $\theta_G=0$}
\label{sec:caseA}

\textbf{Claim to test:} ``If public capital isn't productive ($\theta_G=0$),
then it shouldn't matter whether government spending is transfers or public
investment --- both are just money moving around.''

\textbf{Why this is wrong.} Transfers $G_T$ are \emph{not} resource-using:
they cancel out of the resource constraint $Y=C+I+I^G$ entirely. Public
investment $I^G$ \emph{is} resource-using, with or without $\theta_G>0$ ---
it is a real claim on output that competes with $C$ and $I$ for the same
$Y$. So shifting the composition of a fixed total $G/Y$ from transfers
toward public investment is economically identical to raising ordinary,
unproductive government purchases: a negative wealth effect that makes the
household work more, even though public capital does nothing for
productivity.

\begin{casebox}
\textbf{Numerical example.} $\theta_N=0.58$, $\delta=10\%$, $r=6.5\%$,
$G/Y=20\%$, $\theta_G=0$, lump-sum financing, $\theta_L$ calibrated once at
$I^G/G=0\%$ so that $N=20.0\%$ there. This is exactly what \textbf{Panel 1
(Cumulative Effects)}, Section~\ref{sec:panels}, shows if only the
composition slider (item 5) is moved away from the benchmark's $I^G/G=0\%$.

\begin{center}
\begin{tabular}{ccccc}
\toprule
$I^G/G$ share & $N$ & $Y$ & $C$ & $I^G$ \\
\midrule
$0\%$ (all transfers)   & $20.0\%$ & $0.3934$ & $0.2933$ & $0.0000$ \\
$50\%$                  & $22.4\%$ & $0.4407$ & $0.2845$ & $0.0441$ \\
$100\%$ (all investment) & $25.5\%$ & $0.5009$ & $0.2732$ & $0.1002$ \\
\bottomrule
\end{tabular}
\end{center}
\end{casebox}

$N$, $Y$, and $I$ all rise as the investment share rises, purely from the
resource-cost channel --- \emph{before} any productivity effect is even
switched on. This is exactly why, at $\theta_G=0$, \textbf{Panel 2} still
reports a nonzero headline multiplier $\Delta Y/\Delta G$ for an ordinary,
unproductive spending increase ($\approx 1.1$--$1.2$ under lump-sum
financing) --- $\theta_G=0$ turns off the \emph{productivity} channel, not
the resource-cost channel that drives this table.

\begin{notebox}
\textbf{The one truly neutral corner case.} If (and only if) financing is
lump-sum \emph{and} the composition is already $100\%$ transfers
($I^G/G=0$), then changing the \emph{level} of $G/Y$ is a pure wash: a
lump-sum tax funding an equal-sized lump-sum transfer leaves $Y,C,I,K,N$
completely flat, since there is no $(1-\tau)$ wedge to distort anything and
no resource-using spending to crowd anything out. This is the \emph{only}
combination of settings under which government spending has literally zero
real effect.
\end{notebox}

\section{Case B: the financing rule moves the economy even at $100\%$ transfers}
\label{sec:caseB}

\textbf{Claim to test:} ``If every dollar of tax revenue is handed straight
back to households as a transfer, it shouldn't matter whether it was
collected as a lump-sum poll tax or an income tax --- the household ends up
with the same net income either way.''

\textbf{Why this is wrong.} The claim confuses the \emph{level} of net
transfers with the \emph{marginal} incentive to work and save. Under income
tax financing, the household's capital-Euler condition contains a
$(1-\tau)$ wedge on the \emph{pre-tax} marginal product of capital:
$(1-\tau)\,\theta_K A \kappa^{\theta_K-1}=r+\delta$, instead of
$\theta_K A \kappa^{\theta_K-1}=r+\delta$. To earn the same after-tax return
$r$, the household needs a \emph{higher} pre-tax marginal product, which
means holding \emph{less} capital per worker $\kappa$ in equilibrium --- a
real, permanent effect on the capital stock that has nothing to do with
where the tax revenue eventually lands. The same logic applies to the
labor-leisure condition.

\begin{casebox}
\textbf{Numerical example.} $\theta_N=0.58$, $\delta=10\%$, $r=6.5\%$,
$G/Y=20\%$, all of it transfers ($I^G/G=0\%$), $\theta_G=0$. $\theta_L$ is
calibrated \emph{separately} under each financing rule so that $N=20.0\%$ in
both cases --- this isolates the capital-stock channel from any mechanical
change in hours. \textbf{Panel 1 (Cumulative Effects)}, Section~\ref{sec:panels},
shows this same financing switch, but holds $\theta_L$ fixed at the
benchmark's calibrated value throughout, so $N$ moves too under income tax
(it does not recalibrate back to $20.0\%$) --- the mechanism is identical,
only this table's presentation isolates it more cleanly.

\begin{center}
\begin{tabular}{cccccc}
\toprule
Financing & $N$ & $Y$ & $C$ & $K$ & $w$ \\
\midrule
Lump-sum   & $20.0\%$ & $0.3934$ & $0.2933$ & $1.0014$ & $1.1409$ \\
Income tax & $20.0\%$ & $0.3347$ & $0.2666$ & $0.6816$ & $0.9707$ \\
\bottomrule
\end{tabular}
\end{center}

Switching from lump-sum to income tax financing, with $G/Y$, composition,
and $\theta_G$ all unchanged, cuts $Y$ by $\approx 15\%$ and $K$ by
$\approx 32\%$. $N$ is unchanged only because $\theta_L$ was
\emph{recalibrated} to hit the same target hours under each rule --- if
$\theta_L$ is instead held fixed at one rule's calibrated value (as the app
does when comparing two \emph{policies} under the \emph{same} financing
rule), $N$ falls too under income tax, since the labor-leisure wedge
directly discourages work at the margin.
\end{casebox}

\begin{notebox}
\textbf{Takeaway.} There is no ``neutral financing switch'' the way there
is a neutral composition/level corner case in Case A. The distortion lives
on the household's first-order conditions --- the \emph{margin} --- not on
the government's balance sheet, so it cannot be undone by any choice of
what the revenue funds.
\end{notebox}

\section{Case C: four steady-state variables always move by the identical percentage}
\label{sec:caseC}

\textbf{Claim to test:} ``If two variables show the exact same percent
change in a comparison table, that must be a coincidence or a bug.''

\textbf{Why this is wrong.} Three of the model's accounting identities are
exact, not approximate, and they chain together to force four variables ---
total government spending $G$, public investment $I^G$, transfers $G_T$,
and public capital $K^G$ --- to move by \emph{exactly} the same percentage
whenever the composition shares and $\delta$ are held fixed between the two
points being compared:
\eq{ I^G = f_{I^G}\cdot G, \qquad G_T = f_{G_T}\cdot G \qquad \Longrightarrow \qquad \%\Delta I^G = \%\Delta G_T = \%\Delta G }
\eq{ K^G = \frac{I^G}{\delta} \qquad \Longrightarrow \qquad \%\Delta K^G = \%\Delta I^G }

Chaining these: $\%\Delta K^G=\%\Delta I^G=\%\Delta G=\%\Delta G_T$, always,
regardless of what is \emph{driving} the change in $G$ (a change in $\theta_G$,
$r$, financing, or $G/Y$ itself). The \emph{levels} of these four variables
can differ enormously --- $K^G$ is typically many times larger than $I^G$,
since $K^G=I^G/\delta$ with $\delta$ small --- but the \emph{ratio} between
any two comparison points is forced to be identical by these three
equations, so their percent changes print identically in any comparison
table. This is a feature of the model's bookkeeping, not evidence of an
error.

\section{Case D: the model can lose its saddle path at extreme income-tax rates}
\label{sec:caseD}

\textbf{Claim to test:} ``The app should be able to solve \emph{any}
combination of sliders --- if it throws an error, that's a bug to file.''

\textbf{Why this is (usually) wrong.} The transition path of
Section~\ref{sec:panels} requires the reduced-form matrix $M$ (built from
$\Phi_{kk},\Phi_{kc},\text{coef}_k,\text{coef}_c$) to have exactly one
eigenvalue inside the unit circle and one outside --- the textbook
saddle-path (determinacy) condition for one predetermined and one
jump variable. At sufficiently extreme income-tax rates, this condition
can genuinely fail: \emph{both} eigenvalues end up outside the unit
circle, and \textbf{no bounded perfect-foresight path exists} for the
linearized model. This is not a numerical bug --- it is the honest answer
that this particular calibration has no well-defined local equilibrium
dynamics.

\textbf{Mechanism.} Under income-tax financing, a higher $G/Y$ raises
$\tau$, which lowers the steady-state investment share $s_I$ (more of $Y$
is taxed away, financing government spending instead of private
investment). The capital-accumulation coefficient
$\Phi_{kk}=(1-\delta)+\delta\,y_k/s_I$ \emph{rises} as $s_I$ falls --- a
smaller investment share means the same percentage swing in output
translates into a larger swing in the capital/output ratio, amplifying
the feedback loop between capital and consumption until the ``stable''
root of $M$ is pushed up through $1$.

\begin{casebox}
\textbf{Numerical example.} $\theta_N=0.58$, $\delta=10\%$, $r=6.5\%$,
$\theta_G=0$, all transfers ($I^G/G=0\%$), income-tax financing. The
stable eigenvalue of $M$ as $G/Y$ rises:
\begin{center}
\begin{tabular}{ccccc}
\toprule
$G/Y$ & $\tau$ & $N$ & $s_I$ & Stable eigenvalue of $M$ \\
\midrule
$30\%$ & $0.30$ & $17.5\%$ & $0.178$ & $0.924$ \\
$35\%$ & $0.35$ & $16.2\%$ & $0.165$ & $0.975$ \\
$36\%$ & $0.36$ & $16.0\%$ & $0.163$ & $0.989$ \\
$36.5\%$ & $0.365$ & $15.9\%$ & $0.162$ & $0.997$ \\
$37\%$ & $0.37$ & $15.7\%$ & $0.160$ & $1.006$ \textbf{(fails)} \\
\bottomrule
\end{tabular}
\end{center}
Between $G/Y=36\%$ and $37\%$, the stable root crosses $1$ and the
saddle-path condition fails --- with \emph{zero} composition (100\%
transfers), this is the app's own breaking point. Raising the
public-investment share pushes the threshold higher, since diverting
spending into $I^G$ (rather than transfers) partially offsets the
resource-cost squeeze on $s_I$'s dynamics: the same failure occurs at
$G/Y\approx 39\%$ at $I^G/G=25\%$, $\approx 43\%$ at $50\%$,
$\approx 48\%$ at $75\%$, and $\approx 58\%$ at $100\%$.
\end{casebox}

\begin{notebox}
\textbf{What the app does about it.} Rather than crashing, the app falls
back to the two well-defined steady states (Section~\ref{sec:steadystate}
never requires saddle-path stability --- only the \emph{transition}
between them does) and shows a scoped warning. The long-run comparisons
and long-run multiplier ($\Delta Y/\Delta G$ from the two steady states
directly) remain exact and valid; only the transition-path chart, the
impact multiplier, and the duration-sensitivity replication are skipped,
since they need a well-defined path from one steady state to the other.
There is no single ``safe'' $G/Y$ cutoff to quote in class --- the
threshold depends on financing \emph{and} composition jointly, as the
table above shows.
\end{notebox}

\section{Case E: ``temporary'' does not mean ``back to the benchmark''}
\label{sec:caseE}

\textbf{Claim to test:} ``If I set Duration to Temporary and shrink $G/Y$
back to $20\%$ after a few years, Panel 1's chart should return to
exactly the benchmark, since that's where $G/Y$ started.''

\textbf{Why this is wrong.} Panel 1 compares the sidebar's current
settings against the paper's \emph{fixed} benchmark (lump-sum, $\theta_G=0$,
all transfers) --- but financing, composition, and $\theta_G$ are always
treated as a \emph{permanent} regime switch (Section~\ref{sec:panels};
there is no mechanism for a tax \emph{rule} to revert). So when $G/Y$
reverts to $20\%$ under a Temporary shock, the economy does \emph{not}
return to the benchmark --- it settles at a \emph{third} steady state:
the \emph{current} regime (e.g., income-tax financing) evaluated at the
\emph{benchmark's} $20\%$ spending share. Call this $\text{ss}_{\text{reverted}}$.

\begin{casebox}
\textbf{Numerical example.} $\theta_N=0.58$, $\delta=10\%$, $r=6.5\%$,
$\theta_G=0$, all transfers, income-tax financing, sidebar $G/Y$ moved to
$30\%$.
\begin{center}
\begin{tabular}{lccc}
\toprule
Steady state & $G/Y$ & Financing & $Y$ \\
\midrule
Benchmark ($\text{ss}_{\text{from}}$) & $20\%$ & Lump-sum & $0.3934$ \\
Current ($\text{ss}_{\text{to}}$) & $30\%$ & Income tax & $0.2081$ \\
Reverted ($\text{ss}_{\text{reverted}}$) & $20\%$ & Income tax & $0.2639$ \\
\bottomrule
\end{tabular}
\end{center}
Even after $G/Y$ fully reverts to the benchmark's own $20\%$ share,
output settles $33\%$ \emph{below} the benchmark ($0.2639$ vs.\ $0.3934$)
--- not because of spending, but because income-tax financing is still in
effect. Reverting the \emph{level} of spending did nothing to undo the
\emph{financing rule}.
\end{casebox}

\begin{notebox}
\textbf{A numerical-methods postscript.} Naively modeling ``temporary''
as $G$ reverting to a nonzero constant \emph{forever} (rather than
exactly $\text{ss}_{\text{reverted}}$'s own level, i.e.\ $\hat g_t=0$
relative to that reference) breaks the saddle-path solver's implicit
terminal condition and produces an explosive, diverging path instead of
convergence to $\text{ss}_{\text{reverted}}$. The fix is to linearize
around $\text{ss}_{\text{reverted}}$ itself for the post-reversion phase
--- exactly mirroring how a temporary $\Delta G/Y$ shock in Panel 2
already linearizes around the \emph{old} steady state rather than the
new one (Section~\ref{sec:panels}). The lesson generalizes: in a
perfect-foresight model, a ``temporary'' shock is only well-posed once
you have correctly identified \emph{where it is temporary from} ---
which reference point the economy actually returns to.
\end{notebox}

\section{Case F: the real interest rate changes both the economy's scale and the fiscal multiplier}
\label{sec:caseF}

\textbf{Claim to test:} ``The real interest rate $r$ (sidebar item 3) is
just the household's discount rate --- it should shift the \emph{level}
of the economy up or down, but it shouldn't have anything to do with
\emph{fiscal policy}, since $r$ and $G/Y$ are two completely separate
sliders.''

\textbf{Why this is wrong.} $r$ enters the capital-Euler condition
directly, $\kappa=\big(\text{tax\_factor}\cdot\theta_K A/(r+\delta)\big)^{1/(1-\theta_K)}$
(Section~\ref{sec:steadystate}): a higher $r$ raises the return households
demand on capital, so the steady-state capital/labor ratio $\kappa$ ---
and with it $Y$, $K$, $C$, and $w$ --- fall monotonically as $r$ rises,
\emph{independent of fiscal policy entirely}. That part matches the
claim. But $\kappa$ also determines the investment share
$s_I=\delta\kappa/(Y/N)$, which is one of the ``great ratios'' feeding
directly into the labor-leisure condition and, through it, into
\emph{every} fiscal-policy multiplier in Panels 2--3. $r$ and $G/Y$ are
separate sliders, but they are \emph{not} separate channels in the
model's mechanics.

\begin{casebox}
\textbf{Numerical example --- levels.} $\theta_N=0.58$, $\delta=10\%$,
$G/Y=20\%$, all transfers, lump-sum financing, $\theta_G=0$, $\theta_L$
calibrated once at $r=6.5\%$ and held fixed (Case B's convention, so $N$
is free to move honestly):
\begin{center}
\begin{tabular}{cccccc}
\toprule
$r$ & $\kappa$ & $Y$ & $K$ & $N$ & $s_I$ \\
\midrule
$1.0\%$  & $10.074$ & $0.6112$ & $2.3335$ & $23.2\%$ & $0.382$ \\
$6.5\%$  & $5.007$  & $0.3934$ & $1.0014$ & $20.0\%$ & $0.255$ \\
$12.0\%$ & $3.049$  & $0.2990$ & $0.5709$ & $18.7\%$ & $0.191$ \\
\bottomrule
\end{tabular}
\end{center}

\textbf{Numerical example --- the long-run multiplier $\Delta Y/\Delta G$}
(a marginal, permanent $\Delta G/Y=1$pp increase, $\theta_L$ recalibrated
at each $r$ so $N=20\%$ at the pre-shock baseline, matching Panel 2's own
convention):
\begin{center}
\begin{tabular}{ccc}
\toprule
$r$ & Lump-sum, all investment & Income tax, all transfers \\
\midrule
$1.0\%$  & $+1.38$ & $-4.57$ \\
$6.5\%$  & $+1.13$ & $-3.92$ \\
$12.0\%$ & $+1.04$ & $-3.68$ \\
\bottomrule
\end{tabular}
\end{center}
\end{casebox}

Both multipliers shrink in \emph{magnitude} as $r$ rises --- the
lump-sum/investment multiplier drifts down toward $+1$ (the point where
the extra output exactly equals the extra spending, no supply-side
amplification left), and the income-tax multiplier's negative effect
shrinks too. The mechanism is the falling investment share $s_I$: at a
high $r$, the steady state already sits in a capital-scarce, low-$\kappa$
region where investment is a smaller share of output to begin with, so
the same percentage shift in labor supply or capital demand has less
capital to act on and leverages a smaller output response.

\begin{notebox}
\textbf{Classroom connection.} This is the model's version of a very
current policy debate: fiscal multipliers are widely believed to be
\emph{larger} in low-interest-rate environments (part of the
post-2008/post-2020 case for stimulus when policy rates were near zero)
and smaller when rates are high. The table above shows this is not
just a New Keynesian, zero-lower-bound phenomenon --- it shows up even in
this flexible-price, perfect-foresight real business cycle model, purely
through the investment-share channel above.
\end{notebox}

\section{Case G: longer-lasting temporary shocks have \emph{bigger} impact multipliers}
\label{sec:caseG}

\textbf{Claim to test:} ``A temporary spending increase should have a
\emph{smaller} effect on impact than a permanent one --- if households
know the spending burden will end soon, standard consumption-smoothing
(Barro--Hall) logic says they should barely react today, and react more
and more as the shock is made longer-lasting and eventually permanent, in
a smoothly diminishing way that a short 1-year shock should have almost
\emph{no} impact-year effect at all, essentially zero.''

\textbf{Why the direction is right but the intuition is backwards.} It is
true that a longer-duration shock produces a bigger impact-year response
--- but the reason is the \emph{opposite} of simple consumption-smoothing.
Consumption smoothing is about \emph{demand}: with a short-lived burden,
households can borrow/save their way through it, so consumption barely
moves. But the app's multiplier is driven by the \emph{labor-supply}
response to a negative wealth effect (Case A), and the household's
Euler equation (Section~\ref{sec:steadystate}) makes today's choices
forward-looking over the \emph{entire announced path} of $g_t$: a
government-spending increase that lasts $30$ years is a much
\emph{larger total burden}, in present-value terms, than the same-sized
increase lasting $1$ year. A bigger total burden means a bigger
immediate negative wealth effect, so households work \emph{more}, right
away --- even though the shock hasn't actually gotten bigger \emph{per
year}, only \emph{longer}.

\begin{casebox}
\textbf{Numerical example.} $\theta_N=0.58$, $\delta=10\%$, $r=6.5\%$,
$\theta_G=0$, lump-sum financing, all public investment ($I^G/G=100\%$,
so the resource-cost channel of Case A is active), a marginal, temporary
$\Delta G/Y$ increase of varying duration:
\begin{center}
\begin{tabular}{cc}
\toprule
Duration & Impact (year-0) multiplier $\Delta Y_0/\Delta G$ \\
\midrule
$1$ year & $0.18$ \\
$2$ years & $0.32$ \\
$4$ years & $0.52$ \\
$8$ years & $0.73$ \\
$15$ years & $0.84$ \\
$30$ years & $0.87$ \\
Permanent ($\infty$) & $0.89$ \\
\bottomrule
\end{tabular}
\end{center}
The long-run multiplier is $1.13$ throughout this table --- it is a
comparative-statics number defined purely by the \emph{two steady
states} (Section~\ref{sec:steadystate}), and does not depend on how the
economy gets from one to the other, so it is the same whether or not the
shock is ever actually made permanent. The impact multiplier, by
contrast, rises monotonically with duration and \textbf{asymptotes to,
but never exceeds, the permanent-shock impact multiplier} ($0.89$): a
finite-duration shock is always, in present-value terms, a smaller
total burden than the permanent one.
\end{casebox}

\begin{notebox}
\textbf{Where to see this live.} This is exactly the ``How much does the
\emph{duration} of a temporary shock matter?'' expander under Panel 2
(Section~\ref{sec:panels}), which plots the impact multiplier against
duration for whatever financing/composition/$\theta_G$/G/Y the sidebar
currently has, with the permanent-shock value drawn as a dashed
reference line. Baxter \& King make exactly this point in their
Section~IV, calling it the opposite of the ``Barro--Hall'' intuition that
temporary fiscal shocks should matter less than permanent ones. It is
the model's own version of ``forward guidance'' in modern monetary
policy language: how long a policy is \emph{expected to last} changes
its effect today, independent of its per-period size.
\end{notebox}

\section{Summary table}

\begin{center}
\begin{tabular}{p{0.32\textwidth}p{0.6\textwidth}}
\toprule
Question & Answer \\
\midrule
Does raising $I^G/G$ at $\theta_G=0$ do anything? &
\textbf{Yes} --- it acts like ordinary unproductive spending (Case A). \\[4pt]
Is there \emph{any} setting where changing $G/Y$ does nothing at all? &
\textbf{Yes, exactly one:} lump-sum financing \emph{and} $100\%$ transfers
(Case A, notebox). \\[4pt]
Does switching financing alone matter if all revenue is transferred back? &
\textbf{Yes} --- the $(1-\tau)$ wedge is on the margin, not on net income
(Case B). \\[4pt]
Why do $G$, $I^G$, $G_T$, $K^G$ often show identical \% change? &
Because $I^G=f_{I^G}G$, $G_T=f_{G_T}G$, and $K^G=I^G/\delta$ are exact
identities whenever composition shares and $\delta$ are unchanged (Case C). \\[4pt]
Why does the app sometimes refuse to compute a transition path? &
Extreme income-tax rates (with a low investment share) can push the
model's saddle path to lose determinacy --- a genuine breakdown of the
linearized dynamics, not a bug (Case D). \\[4pt]
Does a ``temporary'' G/Y increase return the economy to the benchmark? &
\textbf{No} --- only spending reverts; financing/composition/$\theta_G$
don't, so the economy settles at a third steady state, not the original
benchmark (Case E). \\[4pt]
Does the real interest rate $r$ affect fiscal multipliers, or only levels? &
\textbf{Both} --- $r$ sets $\kappa$ and hence the investment share $s_I$,
which feeds directly into every multiplier; higher $r$ shrinks multiplier
magnitude (Case F). \\[4pt]
Should a longer-lasting temporary shock have a \emph{smaller} impact
effect (consumption-smoothing intuition)? &
\textbf{No} --- the \emph{opposite}: longer duration means a bigger
present-value burden, so the labor-supply response (and hence the
impact multiplier) rises monotonically with duration, asymptoting to
the permanent-shock value (Case G). \\
\bottomrule
\end{tabular}
\end{center}

\appendix
\part{Appendix: full derivations}

This appendix derives, from scratch, every equation asserted without proof
in Part I: the two household first-order conditions (labor-leisure and
Euler) for both financing rules, and the closed-form steady-state
equations built on top of them. Nothing here is new economics --- it is
the algebra that Part I skipped over.

\section{The firm's problem (why $w_t=\mathrm{MPL}_t$ and $r_t=\mathrm{MPK}_t$)}
\label{app:firm}

Before touching the household, it is worth pinning down where the prices
$w_t$ and $r_t$ in the household's budget constraint actually come from.
A representative, perfectly competitive firm rents private capital and
labor from households each period and uses the publicly provided capital
stock $K^G_t$ for free (it is non-rival and supplied by the government, not
purchased in a market). The firm solves, period by period,
\eq{ \max_{K_t,\,N_t}\ \ A\,K_t^{\theta_K}(K^G_t)^{\theta_G}N_t^{\theta_N} \;-\; w_tN_t \;-\; r_tK_t. }
Since $K^G_t$ is taken as given (external to the firm's own choice), the
first-order conditions are the ordinary marginal-product conditions:
\eq{ \frac{\partial Y_t}{\partial N_t} = w_t \;\Longrightarrow\; w_t=\theta_N\,\frac{Y_t}{N_t} \equiv \mathrm{MPL}_t, }
\eq{ \frac{\partial Y_t}{\partial K_t} = r_t \;\Longrightarrow\; r_t=\theta_K\,\frac{Y_t}{K_t} \equiv \mathrm{MPK}_t. }
Zero profit follows automatically from constant returns to scale in
$(K_t,N_t)$ (Euler's homogeneous-function theorem): at these prices the
firm earns exactly $\theta_K Y_t+\theta_N Y_t=Y_t$, paying out all of
output to the two rented factors. This is why $r_t$ (the household's
capital income rate) and $\mathrm{MPK}_t$ can be used interchangeably
below --- they are the \emph{same number} in every period, by the firm's
optimality condition, not by assumption.

\section{Setting up the household's problem}
\label{app:setup}

The household chooses sequences $\{C_t,N_t,K_{t+1}\}_{t=0}^{\infty}$ to
maximize discounted lifetime utility,
\eq{ \max\ \sum_{t=0}^{\infty}\beta^t\Big[\ln C_t + \theta_L\ln(1-N_t)\Big], \qquad \beta=\frac{1}{1+r}, }
subject to one budget constraint per period and the capital accumulation
identity $K_{t+1}=(1-\delta)K_t+I_t$, i.e.\ $I_t=K_{t+1}-(1-\delta)K_t$.
Substituting this into the two budget constraints of
Section~5 gives, respectively,
\eq{ \text{Income tax:}\qquad C_t + K_{t+1}-(1-\delta)K_t = (1-\tau_t)\big(w_tN_t+r_tK_t\big)+G_{T,t}, \label{app:bc-it} }
\eq{ \text{Lump-sum:}\qquad C_t + K_{t+1}-(1-\delta)K_t = w_tN_t+r_tK_t-\tau_tY_t+G_{T,t}. \label{app:bc-ls} }
There is no expectation operator anywhere in this problem: the model has
no stochastic element, so the household is choosing a sequence under
\textbf{perfect foresight} --- it knows $w_t,r_t,\tau_t,G_{T,t}$ (and hence
$Y_t$) at every future date along the equilibrium path, given the policy
being analyzed. Two modeling conventions are worth flagging explicitly:
\begin{enumerate}[label=(\roman*)]
\item Under \emph{income tax} financing, the wedge $(1-\tau_t)$ multiplies
  the household's \emph{own} factor income $w_tN_t+r_tK_t$ --- it is a
  genuine proportional tax on the household's choices, so it will show up
  when we differentiate with respect to $N_t$ and $K_{t+1}$ below.
\item Under \emph{lump-sum} financing, the household's tax bill is written
  as $\tau_tY_t$, using \emph{aggregate} output $Y_t$, not the household's
  own income. The (representative) household takes $\tau_tY_t$ as a
  parameter it cannot affect with its own $N_t$ or $K_t$ choice --- exactly
  what makes a lump-sum tax lump-sum. This is why $\tau_t$ never appears
  in the lump-sum first-order conditions below, even though $\tau_tY_t>0$
  in general.
\end{enumerate}

\section{Deriving the labor-leisure and Euler conditions}
\label{app:focs}

Attach a Lagrange multiplier $\lambda_t\ge 0$ (period-$t$'s marginal
utility of an extra unit of income) to the period-$t$ budget constraint and
form the Lagrangian. Under \textbf{income tax} financing,
\eq{
\mathcal{L} = \sum_{t=0}^{\infty}\beta^t\Big\{
  \ln C_t + \theta_L\ln(1-N_t)
  + \lambda_t\Big[(1-\tau_t)(w_tN_t+r_tK_t)+G_{T,t}-C_t-K_{t+1}+(1-\delta)K_t\Big]
\Big\}.
}

\paragraph{Step 1: the consumption FOC.} Differentiating with respect to
$C_t$ (which appears only in period $t$'s bracket) gives
\eq{ \frac{\partial\mathcal{L}}{\partial C_t} = \beta^t\left[\frac{1}{C_t}-\lambda_t\right]=0 \;\Longrightarrow\; \lambda_t=\frac{1}{C_t}. \label{app:lambda} }
This is the standard result that the multiplier on the budget constraint
equals the marginal utility of consumption --- an extra unit of income is
worth exactly what an extra unit of consumption is worth, since the
household can always convert one into the other one-for-one within a
period.

\paragraph{Step 2: the labor-leisure FOC.} Differentiating with respect to
$N_t$ (also period-$t$-only) gives
\eq{
\frac{\partial\mathcal{L}}{\partial N_t}
= \beta^t\left[-\frac{\theta_L}{1-N_t}+\lambda_t(1-\tau_t)w_t\right]=0.
}
(The $-\theta_L/(1-N_t)$ term has a sign flip built in: $\partial
\ln(1-N_t)/\partial N_t=-1/(1-N_t)$.) Solving and substituting
$\lambda_t=1/C_t$ from Step 1,
\eq{ \frac{\theta_L}{1-N_t}=\frac{(1-\tau_t)w_t}{C_t} \;\Longrightarrow\; \boxed{\theta_L\,\frac{C_t}{1-N_t}=(1-\tau_t)\,w_t.} }
This is exactly the labor-leisure condition of Section~6. Economically:
the left side is the marginal rate of substitution between consumption and
leisure (how many units of consumption the household would give up for one
more unit of leisure); the right side is the after-tax wage, the market
price of leisure in consumption units. Optimality equates the two.

\paragraph{Step 3: the Euler (capital) FOC.} $K_{t+1}$ is the one variable
that appears in \emph{two} budget constraints: negatively in period $t$'s
(it is being accumulated) and positively in period $t+1$'s (it earns
capital income and is then reduced by depreciation and reinvestment).
Collecting both appearances,
\eq{
\frac{\partial\mathcal{L}}{\partial K_{t+1}}
= \beta^t(-\lambda_t) + \beta^{t+1}\lambda_{t+1}\Big[(1-\tau_{t+1})r_{t+1}+(1-\delta)\Big] = 0.
}
Dividing through by $\beta^t$ and rearranging,
\eq{ \lambda_t = \beta\,\lambda_{t+1}\Big[1+(1-\tau_{t+1})r_{t+1}-\delta\Big]. }
Substituting $\lambda_t=1/C_t$, $\lambda_{t+1}=1/C_{t+1}$, and
$r_{t+1}=\mathrm{MPK}_{t+1}$ (Appendix~\ref{app:firm}),
\eq{ \boxed{\frac{1}{C_t} = \beta\,\frac{1+(1-\tau_{t+1})\,\mathrm{MPK}_{t+1}-\delta}{C_{t+1}}.} }
This is exactly the income-tax Euler equation of Section~6. The
bracketed term $1+(1-\tau_{t+1})\mathrm{MPK}_{t+1}-\delta$ is the
after-tax gross return on one unit of capital held from $t$ to $t+1$: it
pays $(1-\tau_{t+1})\mathrm{MPK}_{t+1}$ in after-tax rental income and
loses $\delta$ to depreciation, on top of the principal. The household is
indifferent at the margin between consuming a unit today (worth $1/C_t$ in
utility, by Step~1) and investing it to consume the gross return tomorrow
(worth $\beta\times\text{return}\times 1/C_{t+1}$ in discounted utility).

\paragraph{The lump-sum case.} Repeat Steps 2--3 with the lump-sum
Lagrangian,
\eq{
\mathcal{L} = \sum_{t=0}^{\infty}\beta^t\Big\{
  \ln C_t + \theta_L\ln(1-N_t)
  + \lambda_t\Big[w_tN_t+r_tK_t-\tau_tY_t+G_{T,t}-C_t-K_{t+1}+(1-\delta)K_t\Big]
\Big\}.
}
Because $\tau_tY_t$ does not depend on $N_t$ or $K_{t+1}$ (point (ii) of
Appendix~\ref{app:setup}), it simply drops out when differentiating:
\eq{
\frac{\partial\mathcal{L}}{\partial N_t}=\beta^t\left[-\frac{\theta_L}{1-N_t}+\lambda_tw_t\right]=0
\;\Longrightarrow\; \boxed{\theta_L\,\frac{C_t}{1-N_t}=w_t,}
}
\eq{
\frac{\partial\mathcal{L}}{\partial K_{t+1}}=\beta^t(-\lambda_t)+\beta^{t+1}\lambda_{t+1}\big[r_{t+1}+(1-\delta)\big]=0
\;\Longrightarrow\; \boxed{\frac{1}{C_t}=\beta\,\frac{1+\mathrm{MPK}_{t+1}-\delta}{C_{t+1}}.}
}
No $(1-\tau)$ wedge appears in either condition: the household's marginal
trade-offs between consumption/leisure and consumption today/tomorrow are
completely undistorted, even though the household \emph{does} pay
$\tau_tY_t$ in taxes every period. This is precisely the sense in which a
lump-sum tax is non-distorting --- it changes the household's wealth, but
not the \emph{price} of working an extra hour or saving an extra dollar.

\section{Deriving the closed-form steady state}
\label{app:steadystate}

A steady state is a situation where every ``great ratio'' and every level
variable is constant over time: $C_t=C$, $N_t=N$, $K_t=K$, $Y_t=Y$, and so
on, for all $t$. This section derives every equation of
Section~\ref{sec:steadystate} in order, starting from the FOCs of
Appendix~\ref{app:focs}.

\subsection{Step 1: the capital-Euler condition pins down $\kappa=K/N$}

In steady state $C_t=C_{t+1}=C$ and $\mathrm{MPK}_t=\mathrm{MPK}_{t+1}=\mathrm{MPK}$,
so the Euler equation collapses to an equation with no time subscripts at
all. Income tax case:
\eq{ 1 = \beta\big[1+(1-\tau)\mathrm{MPK}-\delta\big] \;\Longrightarrow\; (1-\tau)\mathrm{MPK} = \frac{1}{\beta}-1+\delta. }
Recall $\beta=1/(1+r)$ was \emph{chosen} so that $1/\beta=1+r$ (the
household's steady-state discount rate is calibrated to match a target
real interest rate $r$). Substituting,
\eq{ \boxed{(1-\tau)\,\mathrm{MPK} = r+\delta.} }
The lump-sum case is identical without the $(1-\tau)$ factor:
$\boxed{\mathrm{MPK}=r+\delta.}$ These are the two boxed conditions quoted
without proof at the start of Section~\ref{sec:steadystate}.

Now expand $\mathrm{MPK}=\theta_K Y/K$ using the technology equation, and
substitute $K=\kappa N$ (definition of $\kappa$):
\eq{
\mathrm{MPK} = \theta_K\,\frac{A K^{\theta_K}(K^G)^{\theta_G}N^{\theta_N}}{K}
= \theta_K A\,K^{\theta_K-1}(K^G)^{\theta_G}N^{\theta_N}
= \theta_K A\,\kappa^{\theta_K-1}N^{\theta_K-1}(K^G)^{\theta_G}N^{\theta_N}.
}
Since $\theta_K+\theta_N=1$ (constant returns to scale in $(K,N)$ alone,
Section~4), the two powers of $N$ combine to $N^{\theta_K-1+\theta_N}=N^0=1$,
leaving
\eq{ \mathrm{MPK} = \theta_K A\,\kappa^{\theta_K-1}\,(K^G)^{\theta_G}. \label{app:mpk-raw} }
Notice this uses the \emph{raw} level $K^G$, not $K^G/N$: no $N$ survives
the cancellation above. To match the per-worker form used in
Section~\ref{sec:steadystate} and in the app, multiply and divide by
$N^{\theta_G}$:
\eq{ (K^G)^{\theta_G} = \left(\frac{K^G}{N}\right)^{\theta_G}N^{\theta_G} \;\Longrightarrow\; \mathrm{MPK} = \theta_K A\,\left(\frac{K^G}{N}\right)^{\theta_G}N^{\theta_G}\,\kappa^{\theta_K-1}. }
Substituting into the boxed conditions above reproduces exactly
\eq{ \text{Income tax:}\quad (1-\tau)\,\theta_K\,A\,(K^G/N)^{\theta_G}\,N^{\theta_G}\,\kappa^{\theta_K-1}=r+\delta, }
\eq{ \text{Lump-sum:}\quad \theta_K\,A\,(K^G/N)^{\theta_G}\,N^{\theta_G}\,\kappa^{\theta_K-1}=r+\delta. }
\begin{notebox}
\textbf{Why bother with the $(K^G/N)^{\theta_G}N^{\theta_G}$ split instead
of just writing $(K^G)^{\theta_G}$?} Because in the fixed-point iteration
the app actually solves numerically, $K^G$ is not known in advance when
$\theta_G>0$ --- it depends on output through $K^G=I^G/\delta$, which
depends on $\kappa$, which is what we are solving for. Working in
per-worker terms ($K^G/N$) lets the iteration guess and update a
\emph{scale-free ratio} first, and then multiply back in the explicit
$N^{\theta_G}$ correction, rather than guessing the level $K^G$ directly.
When $\theta_G=0$, $N^{\theta_G}=1$ and this distinction is invisible; it
only matters once public capital is productive.
\end{notebox}

When $\theta_G=0$, equation~\eqref{app:mpk-raw} has no $K^G$ dependence at
all, and can be solved for $\kappa$ directly:
\eq{ \theta_K A\,\kappa^{\theta_K-1} = \frac{r+\delta}{\text{tax\_factor}} \;\Longrightarrow\; \kappa = \left(\frac{\text{tax\_factor}\cdot\theta_K A}{r+\delta}\right)^{\frac{1}{1-\theta_K}}, }
where $\text{tax\_factor}=(1-\tau)$ under income tax and $1$ under
lump-sum. When $\theta_G>0$, the same formula holds but with $A$ replaced
by $A\,(K^G/N)^{\theta_G}N^{\theta_G}$, which itself depends on $\kappa$
through $K^G=I^G/\delta$ and $Y=\alpha N$ --- hence the nested fixed-point
iteration described in Section~\ref{sec:steadystate}.

\subsection{Step 2: the investment and consumption shares}

With $\kappa$ (and hence $Y/N\equiv\alpha$) in hand, define the
steady-state investment share directly from the capital accumulation
identity $K_{t+1}=(1-\delta)K_t+I_t$: in steady state $K_{t+1}=K_t=K$, so
$I=\delta K$. Dividing by $Y$,
\eq{ s_I \equiv \frac{I}{Y} = \frac{\delta K}{Y} = \frac{\delta\kappa N}{\alpha N} = \frac{\delta\kappa}{\alpha}. }
The resource constraint $Y=C+I+I^G$ (Section~5), divided by $Y$, gives
$1=s_C+s_I+s_{I^G}$, i.e.
\eq{ \boxed{s_C = 1-s_I-s_{I^G}.} }
Transfers $G_T$ never appear in this identity, because they were never in
the resource constraint to begin with (Section~5) --- this is the formal
statement of ``transfers are not resource-using.''

\subsection{Step 3: the labor-leisure condition pins down $N$}

Start from the steady-state labor-leisure condition (income tax case),
\eq{ \theta_L\,\frac{C}{1-N} = (1-\tau)\,w, \qquad w=\theta_N\,\frac{Y}{N}\ \ \text{(Appendix~\ref{app:firm}).} }
Multiply both sides by $(1-N)$ and divide both sides by $Y$, using
$C/Y=s_C$:
\eq{ \theta_L\,s_C = (1-\tau)\,\theta_N\,\frac{1-N}{N}. }
Multiply both sides by $N$:
\eq{ \theta_L\,s_C\,N = (1-\tau)\,\theta_N\,(1-N) = (1-\tau)\theta_N - (1-\tau)\theta_N\,N. }
Collect every term containing $N$ on the left:
\eq{ N\big[\theta_L\,s_C + (1-\tau)\theta_N\big] = (1-\tau)\theta_N. }
Divide through:
\eq{ \boxed{N = \frac{(1-\tau)\,\theta_N}{\theta_L\,s_C+(1-\tau)\,\theta_N}.} }
The lump-sum case follows identically with $\text{tax\_factor}=1$ in place
of $(1-\tau)$ throughout (there is no wedge in the lump-sum labor-leisure
condition to begin with):
\eq{ \boxed{N = \frac{\theta_N}{\theta_L\,s_C+\theta_N}.} }
These reproduce exactly the two boxed $N$-equations of
Section~\ref{sec:steadystate}. Notice the entire derivation only used
\emph{ratios} ($s_C=C/Y$, $\theta_L,\theta_N,\tau$) --- this is why $N$
can be solved before $Y$, $C$, or $K$ are known in levels: once $\kappa$
fixes $\alpha=Y/N$ and $s_C$ (Steps 1--2), $N$ falls out of a single
algebraic equation, and every level variable is then a trivial multiple of
$N$:
\eq{ Y=\alpha N, \qquad K=\kappa N, \qquad C=s_C\,Y, \qquad I=s_I\,Y, \qquad I^G=s_{I^G}\,Y, \qquad G_T=s_{G_T}\,Y. }

\begin{notebox}
\textbf{Where does $\theta_L$ come from?} $\theta_L$ is not observed
directly; it is \emph{calibrated} by inverting the boxed $N$-equation at a
chosen baseline policy and a target hours-worked share $N_{\text{target}}$
(Baxter \& King's Table~1 strategy). Solving the income-tax $N$-equation
for $\theta_L$ instead of $N$:
\eq{ \theta_L = \frac{(1-\tau)\,\theta_N\,(1-N_{\text{target}})}{s_C\,N_{\text{target}}}, }
and analogously without the $(1-\tau)$ factor under lump-sum. This is
exactly \texttt{calibrate\_theta\_L} in the code: it is the \emph{same}
equation as Step 3 above, just solved for the other unknown.
\end{notebox}

\end{document}
